% STATUS: written 2026-10-21, recalled results as is (one environment each, source kind and whole statement)
% Sources read in full: Bresar 1988 (pp. 1003--1006), Bresar 2005 (pp. 411--425, Sections 1--4).
\section{Cusack, Brešar and the semiprime case}\label{sec:jordan-type:semiprime}

Section~\ref{sec:jordan-type:herstein-theorem} gave the theorems of Herstein and Cusack.
Brešar returned to the problem twice. In 1988 he extended the theorem of Herstein to semiprime
rings \cite[p.~1003]{Bresar1988}. In 2005 he studied Jordan derivations into a bimodule \cite[p.~411]{Bresar2005}.
He showed that such a map is a derivation on a part of the ring \cite[p.~411]{Bresar2005}.


The 1988 paper begins with its definitions. All rings are associative \cite[p.~1003]{Bresar1988}.

\begin{definition}[{\cite[Preliminaries, p.~1003]{Bresar1988}}]\label{def:jordan:bresar-1988}
A ring $R$ is said to be $2$-torsion free if $2a=0$ with $a\in R$ implies $a=0$. A ring $R$
is called a prime ring if $aRb=(0)$ implies $a=0$ or $b=0$. A ring $R$ is called a semiprime
ring if $aRa=(0)$ implies $a=0$. An additive mapping $'\colon R\to R$ is called a Jordan
derivation if $(a^{2})'=a'a+aa'$ holds for all $a\in R$.
\end{definition}

Herstein proved that every Jordan derivation on a $2$-torsion free prime ring
is a derivation \cite[p.~1003]{Bresar1988}. Brešar extends this to semiprime rings \cite[p.~1003]{Bresar1988}.

\begin{theorem}[{\cite[Theorem~1, p.~1003]{Bresar1988}}]\label{thm:jordan:bresar-semiprime}
Let $R$ be a $2$-torsion free semiprime ring and let $'\colon R\to R$ be a Jordan
derivation. In this case, $'$ is a derivation.
\end{theorem}

The proof takes several steps. Brešar follows Herstein and writes $a^{b}$ for
$(ab)'-a'b-ab'$ \cite[p.~1004]{Bresar1988}. The first step is an identity, which he records
for completeness \cite[p.~1004]{Bresar1988}.

\begin{theorem}[{\cite[Theorem~3, p.~1004]{Bresar1988}}]\label{thm:jordan:bresar-identity}
Let $R$ be a $2$-torsion free ring and let $'\colon R\to R$ be a Jordan derivation.
In this case for all elements $a,b,x\in R$,
$a^{b}x[a,b]+[a,b]xa^{b}=0$ holds.
\end{theorem}

A second lemma on pairs $a,b$ supplies the rest.

\begin{lemma}[{\cite[Lemma~4, p.~1004]{Bresar1988}}]\label{lem:jordan:bresar-axb}
Let $R$ be a $2$-torsion free semiprime ring and let $a,b\in R$. If $axb+bxa=0$ holds
for all $x\in R$, then $axb=bxa=0$ for all $x\in R$.
\end{lemma}

Brešar proves Lemma~\ref{lem:jordan:bresar-axb} by using the hypothesis three times to
reach $2(axb)y(axb)=0$ \cite[p.~1004]{Bresar1988}. He then combines the lemma with
Theorem~\ref{thm:jordan:bresar-identity} in several steps. These show that each
$a^{b}$ is central with $(a^{b})^{3}=0$, so $a^{b}=0$, since the center of a semiprime
ring has no nonzero nilpotent elements \cite[p.~1005]{Bresar1988}. Two consequences are stated \cite[pp.~1005--1006]{Bresar1988}.

\begin{corollary}[{\cite[Corollary~5, p.~1005]{Bresar1988}}]\label{cor:jordan:bresar-semisimple}
Let $R$ be a $2$-torsion free semisimple ring and let $'\colon R\to R$ be a Jordan
derivation. In this case $'$ is a derivation.
\end{corollary}

The second consequence concerns continuity.

\begin{theorem}[{\cite[Theorem~6, p.~1006]{Bresar1988}}]\label{thm:jordan:bresar-continuous}
Let $A$ be a Banach algebra and let $'\colon A\to A$ be a linear Jordan derivation.
In this case $'$ is continuous.
\end{theorem}

Brešar combines Corollary~\ref{cor:jordan:bresar-semisimple} with the theorem of
Johnson and Sinclair that every linear derivation on a semisimple Banach algebra is
continuous \cite[p.~1006]{Bresar1988}. This answers a question of Sinclair
\cite[p.~1006]{Bresar1988}.


The second paper allows values in a bimodule \cite[p.~411]{Bresar2005}.

\begin{definition}[{\cite[Section~1, p.~411]{Bresar2005}}]\label{def:jordan:bresar-2005}
Let $A$ be a ring and let $M$ be an $A$-bimodule. An additive map $d\colon A\to M$ is called a
Jordan derivation if $d(x^{2})=d(x)x+xd(x)$ and $d(xyx)=d(x)yx+xd(y)x+xyd(x)$ for all
$x,y\in A$.
\end{definition}

When $M$ is $2$-torsion free, the second condition follows from the first
\cite[p.~414]{Bresar2005}. Brešar measures the failure of the product rule by a map
$\{x,y\}=d(xy)-d(x)y-xd(y)$. The map $d$ is a derivation exactly when $\{A,A\}=0$
\cite[p.~414]{Bresar2005}. The map $\{x,y\}$ satisfies the three identities of the next
theorem \cite[p.~414]{Bresar2005}. For a ring $A$ he sets $A^{(0)}=A$ and
$A^{(n+1)}=[A^{(n)},A^{(n)}]$ \cite[p.~412]{Bresar2005}. He writes $R(S)$ and $I(S)$ for the
subring and the ideal of $A$ generated by a subset $S$ \cite[p.~412]{Bresar2005}.

\begin{theorem}[{\cite[Theorem~2.1, p.~413]{Bresar2005}}]\label{thm:jordan:bresar-biadditive}
Let $A$ be a ring and let $M$ be an $A$-bimodule. If $\{\,.\,,\,.\,\}\colon A\times A\to M$
is a biadditive map such that $\{x,x\}=0$, $\{xy,x\}+\{x,y\}x=0$ and
$\{xy,z\}+\{x,y\}z=\{x,yz\}+x\{y,z\}$ for all $x,y,z\in A$, then
$\{R(A^{(2)}),A\}=0$ (and hence $\{I(A^{(3)}),A\}=0$).
\end{theorem}

Applied to the map $\{x,y\}$ above, this gives the main result of the paper.

\begin{theorem}[{\cite[Theorem~3.1, p.~414]{Bresar2005}}]\label{thm:jordan:bresar-restriction}
If $d$ is a Jordan derivation from a ring $A$ into an $A$-bimodule $M$, then
$d(ux)=d(u)x+ud(x)$ for all $x\in A$ and all $u\in R(A^{(2)})$ (and so, in particular,
for all $u\in I(A^{(3)})$).
\end{theorem}

The theorem has an immediate consequence \cite[p.~414]{Bresar2005}.

\begin{corollary}[{\cite[Corollary~3.2, p.~414]{Bresar2005}}]\label{cor:jordan:bresar-big}
If a ring $A$ is such that $R(A^{(2)})=A$ (in particular, if $I(A^{(3)})=A$), then every
Jordan derivation from $A$ into an arbitrary $A$-bimodule is a derivation.
\end{corollary}

Brešar remarks that the theorem gives no information when $A$ is commutative or close
to it \cite[p.~411]{Bresar2005}. For a $2$-torsion free semiprime ring he shows that
$I(A^{(n)})$ is essential exactly when $Z(A)$ contains no nonzero ideal of $A$
\cite[Lemma~3.7, p.~416]{Bresar2005}. This gives the next statement.

\begin{corollary}[{\cite[Corollary~3.8, p.~416]{Bresar2005}}]\label{cor:jordan:bresar-essential}
Let $A$ be a $2$-torsion free semiprime ring and let $d$ be a Jordan derivation from $A$
into an $A$-bimodule $M$.
(i) If $Z(A)$ does not contain nonzero ideals of $A$, then the restriction of $d$ to
some essential ideal of $A$ is a derivation.
(ii) Suppose that $M$ satisfies the following two conditions: (a) for every essential
ideal $E$ of $A$ and every $m\in M$, $Em=0$ implies $m=0$, and (b) $cm=mc$ for all
$c\in Z(A)$ and all $m\in M$. Then $d$ is a derivation.
\end{corollary}

Both conditions of part (ii) hold when $M=A$, so part (ii) extends the theorem of
Cusack \cite[p.~416]{Bresar2005}. Brešar adds that the arguments of Cusack and of a later
proof do not give (ii), since both compute with $d^{2}$ \cite[p.~416]{Bresar2005}.
The next statement concerns the maximal ring of quotients \cite[Corollary~3.9, p.~417]{Bresar2005}.

\begin{corollary}[{\cite[Corollary~3.9, p.~417]{Bresar2005}}]\label{cor:jordan:bresar-quotients}
Let $A$ be a $2$-torsion free semiprime ring and let $Q$ be the maximal left (or right)
ring of quotients of $A$. Then every Jordan derivation from $A$ into $Q$ is a
derivation.
\end{corollary}

For a simple ring Brešar states a definitive conclusion \cite[p.~417]{Bresar2005}.

\begin{corollary}[{\cite[Corollary~3.10, p.~417]{Bresar2005}}]\label{cor:jordan:bresar-simple}
Let $A$ be a noncommutative simple ring with $\operatorname{char}(A)\neq2$. Then every
Jordan derivation from $A$ into any $A$-bimodule is a derivation.
\end{corollary}


Brešar also gives examples of proper Jordan derivations, that is, Jordan derivations
that are not derivations \cite[p.~411]{Bresar2005}. One example changes the second example of
Cusack \cite[\S4.2, p.~419]{Bresar2005}. Let $A$ be the $3$-dimensional algebra over a field $F$ generated by $a$ and
$b$ with $a^{2}=b^{2}=0$ and $ab=-ba\neq0$ \cite[\S4.2, p.~419]{Bresar2005}. Then
$A^{(2)}=0$ and $x\circ y=0$ for all $x,y\in A$, so every linear map from $A$ to itself
is a Jordan derivation \cite[\S4.2, p.~419]{Bresar2005}. Theorem~\ref{thm:jordan:bresar-restriction}
therefore gives no information here. The map $d(\lambda a+\mu b+\nu ab)=\nu a$ is not a
derivation, and it has the following property \cite[\S4.2, p.~419]{Bresar2005}.

\begin{example}[{\cite[Example~4.2, p.~419]{Bresar2005}}]\label{ex:jordan:bresar-three}
There exists a $3$-dimensional noncommutative algebra $A$ admitting a proper Jordan
derivation $d\colon A\to A$ such that for every $u\neq0$ in $A$ there is $x\in A$
satisfying $d(ux)\neq d(u)x+ud(x)$.
\end{example}

Section~\ref{sec:jordan-type:triple-and-generalized} turns to maps that satisfy a
triple version of the product rule.
